% ═══════════════════════════════════════════════════════════════
% MT-01c-ML: Musterlösung Minitest Artikel + Alphabet
% Spanisch Klasse 7 - Sequenz 1: Empezamos
% ═══════════════════════════════════════════════════════════════

\documentclass[11pt, a4paper]{article}

\usepackage[utf8]{inputenc}
\usepackage[T1]{fontenc}
\usepackage[ngerman]{babel}
\usepackage{helvet}
\renewcommand{\familydefault}{\sfdefault}

\usepackage[margin=2cm]{geometry}
\usepackage{enumitem}
\usepackage{tabularx}
\usepackage{booktabs}
\usepackage{xcolor}
\usepackage{fancyhdr}
\usepackage{lastpage}
\usepackage{amssymb}

% FARBEN
\definecolor{spanisch}{HTML}{c41e3a}
\definecolor{grau}{HTML}{888888}
\definecolor{hellgrau}{HTML}{f5f5f5}
\definecolor{loesung}{HTML}{c62828}

\pagestyle{fancy}
\fancyhf{}
\renewcommand{\headrulewidth}{0.4pt}
\renewcommand{\footrulewidth}{0.4pt}

\lhead{Spanisch -- Klasse 7}
\chead{\textbf{MT-01c -- MUSTERLÖSUNG}}
\rhead{\today}

\lfoot{10 Minuten}
\cfoot{}
\rfoot{Seite \thepage\ von \pageref{LastPage}}

\newcommand{\punktebox}[1]{\hfill\fbox{\small \_\_\_/#1}}
\newcommand{\luecke}[1]{\rule{#1}{0.4pt}}
\newcommand{\es}[1]{\textcolor{spanisch}{\textbf{#1}}}
\newcommand{\lsg}[1]{\textcolor{loesung}{\textbf{#1}}}

\begin{document}

% ═══════════════════════════════════════════════════════════════
% KOPFBEREICH
% ═══════════════════════════════════════════════════════════════

\begin{center}
    {\Large\bfseries\textcolor{spanisch}{Minitest: Artikel + Alphabet}}

    \vspace{5pt}
    {\large\textcolor{loesung}{-- MUSTERLÖSUNG --}}
\end{center}

\vspace{10pt}

% ═══════════════════════════════════════════════════════════════
% AUFGABE 1: Singular-Artikel (6 BE)
% ═══════════════════════════════════════════════════════════════

\noindent\textbf{\textcolor{spanisch}{Aufgabe 1 -- Bestimmte Artikel (Singular)}} \punktebox{6}

\vspace{10pt}

\begin{tabularx}{\textwidth}{|l|X|l|X|}
\hline
\lsg{el} & libro (Buch) & \lsg{la} & mesa (Tisch) \\
\hline
\lsg{el} & cuaderno (Heft) & \lsg{la} & silla (Stuhl) \\
\hline
\lsg{el} & chico (Junge) & \lsg{la} & chica (Mädchen) \\
\hline
\end{tabularx}

\vspace{5pt}
\textit{Je 1 BE pro richtigen Artikel. Regel: -o = el (maskulin), -a = la (feminin).}

\vspace{15pt}

% ═══════════════════════════════════════════════════════════════
% AUFGABE 2: Plural-Artikel (6 BE)
% ═══════════════════════════════════════════════════════════════

\noindent\textbf{\textcolor{spanisch}{Aufgabe 2 -- Bestimmte Artikel (Plural)}} \punktebox{6}

\vspace{10pt}

\begin{tabularx}{\textwidth}{|l|X|l|X|}
\hline
\lsg{los} & libros (Bücher) & \lsg{las} & mesas (Tische) \\
\hline
\lsg{los} & chicos (Jungen) & \lsg{las} & sillas (Stühle) \\
\hline
\lsg{los} & profesores (Lehrer) & \lsg{las} & profesoras (Lehrerinnen) \\
\hline
\end{tabularx}

\vspace{5pt}
\textit{Je 1 BE pro richtigen Artikel. Regel: el $\to$ los, la $\to$ las.}

\vspace{15pt}

% ═══════════════════════════════════════════════════════════════
% AUFGABE 3: Genus-Regel (4 BE)
% ═══════════════════════════════════════════════════════════════

\noindent\textbf{\textcolor{spanisch}{Aufgabe 3 -- Genus-Regeln}} \punktebox{4}

\vspace{10pt}

\noindent
a) Nomen auf \textbf{-o} sind meist:

\hspace{1cm} \lsg{$\boxtimes$ maskulin (el)} \hspace{2cm} $\square$ feminin (la) \hfill \textit{(2 BE)}

\vspace{0.8em}

\noindent
b) Nomen auf \textbf{-a} sind meist:

\hspace{1cm} $\square$ maskulin (el) \hspace{2cm} \lsg{$\boxtimes$ feminin (la)} \hfill \textit{(2 BE)}

\vspace{5pt}
\textit{Wichtige Ausnahmen: el día (der Tag), el problema (das Problem), la mano (die Hand).}

\vspace{15pt}

% ═══════════════════════════════════════════════════════════════
% AUFGABE 4: Alphabet Buchstabennamen (6 BE)
% ═══════════════════════════════════════════════════════════════

\noindent\textbf{\textcolor{spanisch}{Aufgabe 4 -- Spanische Buchstabennamen}} \punktebox{6}

\vspace{10pt}

\begin{enumerate}[label=\alph*)]
    \item \textbf{H} = \lsg{hache} \hfill \textit{(2 BE)} \\
    \textit{Besonderheit: Das H ist stumm! hola = [óla]}
    \item \textbf{J} = \lsg{jota} \hfill \textit{(2 BE)} \\
    \textit{Aussprache: wie CH in ``ach''}
    \item \textbf{Ñ} = \lsg{eñe} \hfill \textit{(2 BE)} \\
    \textit{Aussprache: wie NJ in ``Cognac''. Auch akzeptiert: ``ene''}
\end{enumerate}

\vspace{15pt}

% ═══════════════════════════════════════════════════════════════
% AUFGABE 5: Buchstabieren (4 BE)
% ═══════════════════════════════════════════════════════════════

\noindent\textbf{\textcolor{spanisch}{Aufgabe 5 -- Buchstabieren}} \punktebox{4}

\vspace{10pt}

\noindent
\es{HOLA} = \lsg{hache - o - ele - a}

\vspace{5pt}
\textit{Je 1 BE pro richtigen Buchstabennamen. Auch akzeptiert: ohne Bindestriche oder mit Kommas.}

% ═══════════════════════════════════════════════════════════════
% PUNKTEÜBERSICHT
% ═══════════════════════════════════════════════════════════════

\vspace{20pt}
\hrule
\vspace{10pt}

\noindent
\begin{tabularx}{\textwidth}{|X|c|c|c|c|c||c|}
\hline
\textbf{Aufgabe} & 1 & 2 & 3 & 4 & 5 & \textbf{Gesamt} \\
\hline
\textbf{max. Punkte} & 6 & 6 & 4 & 6 & 4 & \textbf{26} \\
\hline
\end{tabularx}

\vspace{15pt}

\noindent\textbf{Notenskala (14-NP):}

\vspace{5pt}

\begin{tabular}{|c|c|c|c|c|c|c|c|c|c|c|c|c|c|c|}
\hline
NP & 14 & 13 & 12 & 11 & 10 & 9 & 8 & 7 & 6 & 5 & 4 & 3 & 2 & 1 \\
\hline
\% & 100 & 96 & 91 & 86 & 80 & 73 & 67 & 60 & 53 & 47 & 40 & 33 & 27 & 20 \\
\hline
BE & 26 & 25 & 24 & 22 & 21 & 19 & 17 & 16 & 14 & 12 & 10 & 9 & 7 & 5 \\
\hline
\end{tabular}

\end{document}
